Large language models (LLMs), renowned for their powerful conversational abilities, are widely recognized as exceptional tools in the field of education, particularly in the context of automated intelligent instruction systems for language learning. In this paper, we propose a scoring system based on LLMs, motivated by their positive impact on text-related scoring tasks. Specifically, the speech encoder first maps the learner’s speech into contextual features. The adapter layer then transforms these features to align with the text embedding in latent space. The assessment task-specific prefix and prompt text are embedded and concatenated with the features generated by the modality adapter layer, enabling the LLMs to predict accuracy and fluency scores. Our experiments demonstrate that the proposed scoring systems achieve competitive results compared to the baselines on the Speechocean762 datasets. Moreover, we also conducted an ablation study to better understand the contributions of the prompt text and training strategy in the proposed scoring system.


% 1000 characters limitation
% A typical fluency scoring system generally relies on an automatic speech recognition (ASR) system to obtain the time stamps of speech segments for further calculation of fluency related handcrafted features or directly modeling of speech fluency as in an end-to-end of approach. This paper proposes a novel ASR-free approach for automatic fluency assessment using self-supervised learning (SSL). Specifically, wav2vec2.0 is used to extract frame-level speech features, followed by K-means clustering to assign a cluster index (pseudo labeling) to each frame. A BLSTM-based scorer is then trained to predict an utterance-level fluency score from the frame-level SSL features along with the corresponding cluster index embeddings. Experimental results conducted on our non-native databases show that the proposed framework significantly improves the performance in the ``open response" scenario compared to previous methods and matches the recently reported performance in the ``read aloud" scenario.

%outperforms baseline systems in the ``open response" scenario, in terms of Pearson correlation coefficients (PCC). In addition, the results imply that the proposed method also preforms well in the ``read aloud" scenario.